../cictikz/src/cictikz/data/tex/cictikz_preamble.tex

% Why a loop of wire is an inductor, told as the three steps the
% equations take. (a) push a changing current through a real loop of
% wire: (1) the current wraps a magnetic field around the wire, so
% flux threads the loop; (2) changing that flux drives a voltage
% around the loop (Faraday); (3) that voltage opposes the change,
% which at the terminals is v = L di/dt, and the ramp/step inset says
% it concretely. (b) wind the same loop N times and both the flux
% made and the flux felt scale with N, hence L grows as N squared.

\begin{tikzpicture}[thick]

  \newcommand{\stepnum}[3]{%
    \draw[fill=white] (#1,#2) circle (0.21);
    \node[anchor=center, scale=0.85] at (#1,#2) {#3};
  }

  % ================= (a) one loop of wire =================
  % a clean loop of wire: an ellipse with a gap at the bottom, its
  % two ends brought out sideways as terminals
  \draw[line width=2.6pt] (-0.62,-1.0) arc (250:-70:1.85 and 1.06);
  \draw[line width=2.4pt] (-0.62,-1.0) -- (-0.62,-1.85) -- (1.5,-1.85);
  \draw[line width=2.4pt] (0.62,-1.0) -- (0.62,-1.5) -- (1.5,-1.5);

  %- (1) the current
  \draw[blue, very thick, ->] (-1.55,1.05) arc (150:30:1.85 and 1.1);
  \node[blue, anchor=south, scale=0.95] at (0,1.65) {$i(t)$, rising};
  \stepnum{-2.35}{0.35}{1}
  \node[anchor=east, align=right, scale=0.85] at (-2.65,0.35)
    {a rising current\\wraps a field};

  %- (2) the flux it makes, threading its own loop
  \foreach \x in {-0.8,0.0,0.8} {
    \draw[red, very thick, ->] (\x,-0.6) -- (\x,0.75);
  }
  \node[red, anchor=west, scale=0.95] at (2.15,0.55) {$\Phi = L\,i$};
  \stepnum{2.35}{-0.2}{2}
  \node[red!60!black, anchor=west, align=left, scale=0.85] at (2.65,-0.2)
    {flux through\\its own loop};

  %- (3) the voltage that changing flux drives, at the terminals
  \draw[armygreen, very thick, ->] (1.5,-1.68) -- (2.4,-1.68);
  \node[armygreen, anchor=west, scale=0.95] at (2.45,-1.68)
    {$v = L\,\dfrac{di}{dt}$};
  \stepnum{-1.3}{-1.85}{3}
  \node[armygreen!60!black, anchor=east, align=right, scale=0.85] at (-1.6,-1.85)
    {and it pushes\\back};

  %- the inset that makes it concrete: ramp the current, get a step
  \begin{scope}[shift={(-1.9,-4.6)}]
    \draw[gray!60, ->] (0,-1.25) -- (3.5,-1.25);
    \draw[blue, very thick] (0.15,0.15) -- (1.1,0.15) -- (2.3,1.15) -- (3.2,1.15);
    \node[blue, anchor=west, scale=0.85] at (3.25,0.85) {$i$};
    \draw[armygreen, very thick] (0.15,-0.95) -- (1.1,-0.95) -- (1.1,-0.35)
      -- (2.3,-0.35) -- (2.3,-0.95) -- (3.2,-0.95);
    \node[armygreen, anchor=west, scale=0.85] at (3.25,-0.55) {$v$};
    \node[anchor=north, scale=0.85, align=center] at (1.7,-1.45)
      {ramp the current, and the loop\\answers with a voltage};
  \end{scope}

  \node[anchor=north, align=center] at (0.4,-7.0)
    {(a) any loop of wire is an inductor};

  % ================= (b) the same loop, N times =================
  \begin{scope}[shift={(9.4,-0.3)}]
    % flux bundle through the bore, behind the winding
    \draw[red, line width=3pt, ->] (-1.1,0.45) -- (5.1,0.45);
    \node[red, anchor=south west] at (3.7,0.6) {$N\Phi$};
    % the winding: one wire, N turns
    \draw[line width=2.4pt] (0,-0.4) -- (-0.8,-0.4);
    \foreach \x in {0,0.8,1.6,2.4,3.2} {
      \draw[line width=2.4pt] (\x,-0.4)
        .. controls (\x-0.25,1.0) and (\x+1.05,1.0) .. (\x+0.8,-0.4);
      \draw[line width=2.4pt, gray!55]
        (\x+0.8,-0.4) .. controls (\x+0.72,-0.75) and (\x+0.1,-0.75) .. (\x+0.02,-0.42);
    }
    \draw[line width=2.4pt] (4.0,-0.4) -- (4.8,-0.4);
    \draw[blue, very thick, ->] (-0.75,-0.7) -- (-0.2,-0.7);
    \node[blue, anchor=north, scale=0.9] at (-0.45,-0.8) {$i$};

    %- why it squares
    \node[anchor=north, align=left, scale=0.9] at (2.0,-1.5)
      {$N$ turns make $N$ times the flux,\\
       and every turn feels all of it:\\
       $L \propto N^2$};
    \node[anchor=north, align=center] at (2.0,-6.7)
      {(b) a coil is the same loop, wound $N$ times};
  \end{scope}

\end{tikzpicture}
\end{document}
